\section{The propers: detailed commentary}
\label{detailed-commentary}
\subsection{Introit: a promise addressed to people in exile}
Jeremias addresses people who have been carried away from Jerusalem to
Babylon. The letter tells them to build houses, plant gardens, raise
families and seek the city's peace. Rival prophets promise a quick reversal.
God's assurance of peace belongs within this long obedience in exile
(Jer 29:1--14). The chant selects the promise, the calling and the return.
It omits much of the letter's instruction and condenses the regathering
clause. Its peaceful future concerns people who must still endure their
present displacement.

Jerome explains the promise by opposing the false prophets' timetable to
God's knowledge. The exiles cannot make a speedy return true by wanting it.
He then pauses over the invitation to pray. God could fulfil his promise
without the captives' invocation, yet urges them to ask so that they may
receive what he has promised. The summons to seek with the whole heart
concerns the persons who receive the gift, as well as the gift they
receive.\footnote{Jerome, \work{Commentary on Jeremias}, book V,
on Jer 29:10--14, PL 24 (Paris, 1865), cols.~893--894; subsequent
context cols.~895--896.}

The appointed psalm verse thanks God for blessing his land and ending
Jacob's captivity. The whole Psalm 84 goes on to ask for renewed life and
salvation. Augustine attends to this combination of completed restoration
and continuing petition. In his Christian exposition, release from captivity
includes forgiveness of the sins that held a person captive. The prayer for
conversion acknowledges God's action in the turning itself.\footnote{Augustine,
\work{Exposition of Psalm 84}, §§2--6; NPNF I.8 prints the heading
Psalm 85. Psalm numbers in this study are Vulgate numbers unless otherwise
stated.} Jeremias's actual exiles and the psalm's Israel remain the
literal subjects. Augustine explains the freedom Christians receive through
Christ by attending to the psalm's own association of release and pardon.

\subsection{Collect: a family dependent on grace}
The Collect names the petitioners God's family. Its relative clause says
that this family relies solely on the hope of heavenly grace. The petition
for constant protection follows from that dependence. Protection must
continue because the family continues to depend on God.

The family asks as a body. Colossians will likewise speak of one body, but
its members still have grievances against one another. Read beside that
Epistle, the Collect permits a precise conclusion: the people commanded
to forgive are themselves people whom God must preserve. Their common
need precedes their examination of a neighbour's faults. This relation
between the prayer and the lesson is especially pertinent when Augustine
asks whether any hearer can be certain of his own perseverance.

\subsection{Epistle: what the forgiven put on}
Colossians 3 contrasts the old life with the new. After commanding the
death of sinful practices and the putting off of anger, malice and lying,
Paul tells God's chosen, holy and beloved people to put on mercy,
kindness, humility, modesty and patience. The clothing image concerns
conduct appropriate to a life already received. Election does not exempt
them from the command to change (Col 3:1--17).

Chrysostom explains mercy as an inward affection, like a parent's concern,
and develops the command to forgive by attending to its example: as the
Lord forgave. A serious injury still falls within this command. The
injured person's innocence in the original quarrel does not release him
from forgiveness. The Lord's action supplies the comparison, rather than
a calculation of who first deserved kindness.\footnote{John Chrysostom,
\work{Homilies on Colossians}, VIII, exposition of Col 3:12--15;
NPNF I.13. His lemma reads ``peace of God'' where the appointed Latin
reads \textit{pax Christi}, peace of Christ.}

Charity is then called the bond of perfection. Chrysostom asks why Paul
calls it a bond. Qualities that appear admirable can come apart when love
is absent: a person may concede a grievance while remaining cold toward
the other. His images are structural and concrete: ligaments join bones,
and the timbers of a house need their fastenings. Love gives the several
virtues their abiding relation to the neighbour.

The end of the Epistle describes instruction and worship. Christ's word
is to dwell abundantly; Christians teach and admonish one another; their
song is addressed to God. Chrysostom applies this abundance to people
responsible for wives and children as well as to religious teachers.
Scripture furnishes enduring comfort under losses because its instruction
remains available after the first trouble has passed.\footnote{Chrysostom,
\work{Homilies on Colossians}, IX, opening exposition of Col 3:16--17.}
Aquinas, on the latter verses, distinguishes teaching, heartfelt praise
and the direction of action to God. Acting in Christ's name includes a
habitual intention that can inform work without a continuously explicit
mental act.\footnote{Thomas Aquinas, \work{Commentary on Colossians},
chapter 3, lectio 3, on verses 16--17, in \work{In omnes Pauli epistolas
commentaria}, II (Dessain, 1857), p.~442. The checked source lacks the
preceding pages of this exposition.} Thanksgiving thus accompanies ordinary
action as well as the assembly's song.

\subsection{Gradual: praise inside a lament}
Psalm 43 first remembers God's gift of the land to the ancestors. Their
sword did not secure it independently of him. The appointed verses
ascribe rescue to God and promise praise in his name. Yet the psalm then
laments defeat, dispersion and reproach. Later it insists that the people
have not forgotten their covenant. Its question concerns the suffering
of those who remain faithful, rather than a simple sequence in which
repentance immediately produces public success.

Augustine hears the martyrs in this psalm, pointing to Paul's use of its
words about sheep for slaughter in Romans 8. He distinguishes past divine
rescue, hoped-for deliverance and present affliction. His exposition of
the appointed rescue also treats the past tense as prophetic certainty:
God can speak of a future gift as accomplished because its fulfilment is
certain. He supports that inference with other prophetic uses of the past
tense and with the psalm's own future praise.\footnote{Augustine,
\work{Exposition of Psalm 43}, §§1,3--8; NPNF I.8 heading Psalm 44.
The translation's verse numbers differ from the Missal's native Ps 43:8--9;
the expository sections identify the argument.}

This Christian reception answers a difficulty the short chant alone
might leave obscure. A person can thank God for rescue while still
enduring persecution. Neither the psalm's own history nor Augustine's
exposition licenses treating every sufferer as someone whom God has
abandoned. Within this Mass the praise can accompany the Gospel's
uncompleted harvest, whose owner has promised a future action.

\subsection{Alleluia and Offertory: two uses of the same cry}
Psalm 129 opens from the depths and proceeds through the acknowledgement
of sin to waiting for the Lord and the hope of Israel's redemption.
The Alleluia selects the cry for hearing and surrounds it with the
acclamation. The Offertory repeats the opening cry at its close.
The repetition belongs to the chant; the biblical English above preserves
the underlying verses rather than translating their adaptation.

Augustine begins with Jonas, whose prayer reached God from within the
fish. Physical barriers do not keep God distant. He then asks what depth
the Christian recognises in himself. A sinner who cries has already begun
to hope; more grievous is the state in which a person no longer thinks
prayer necessary. The psalm's next question, who could stand if God
strictly marked iniquities, prevents the petitioner from claiming a
special innocence.\footnote{Augustine, \work{Exposition of Psalm 129},
§§1--2; NPNF I.8 heading Psalm 130.}

The two liturgical uses allow the same words to be heard before the
Gospel and before the offering. One can infer a relation to the Secret:
those who cry for hearing soon ask for pardon. Augustine's ensuing
exposition identifies the psalm's propitiation with Christ's sacrifice.
His argument receives fuller treatment in the second interpretation
below, where the relation between forgiveness and assured prayer is at
issue.

\subsection{Gospel: the owner's reason for waiting}
Matthew sets this parable among the teaching beside the sea. Jesus later
explains it privately to the disciples in the house. The good sower is
the Son of man; the field is the world; the good seed is the children of
the kingdom; the enemy is the devil; the harvest is the world's
consummation; the reapers are angels (Matt 13:36--43). Those identifications
clarify the parable that the appointed Gospel ends at the barn. They also
explain why its final separation is larger than the servants' proposed
act of weeding.

The servants' question is reasonable. They know that the owner sowed good
seed, and they can now see weeds. He answers their question about origin:
an enemy has done this. Their further request concerns action. His
refusal gives a reason about consequences: uprooting the tares could
uproot the wheat. The story therefore distinguishes recognition of evil
from permission to use a particular remedy against it. The servants need
the owner's command at this point as much as they needed it at sowing.

Chrysostom develops this danger as a prohibition of killing heretics.
Augustine addresses wicked Christians more generally and stresses both
conversion and the risk of human misidentification. Their explanations
are developed in the first interpretation. Both preserve the harvest's
fire and barn. The postponed separation has a stated end; the prohibition
is not a declaration that wheat and tares are identical.

\subsection{Secret: pardon and the direction of the heart}
The Secret calls the offerings a sacrifice of propitiation and asks for
two merciful acts: absolution of offences and guidance of wavering hearts.
The petitioners name their own sins. They have heard the Gospel's account
of an enemy's action, yet their prayer does not confine the need for
mercy to people outside the assembly. They also need direction after
pardon, because a wavering heart can err again.

The prayer's sacrificial language concerns the relation of the offerers
to God. Colossians commands them to forgive each other because the Lord
has forgiven them; it does not replace God's forgiveness with their mutual
agreement. Augustine's account of propitiation in Psalm 129 explains the
Christological ground of a sinner's approach to God. His exposition is
of the psalm, while the Secret itself supplies the actual petition here.

\subsection{Preface: thanksgiving to the triune God}
The common opening gives thanks always and everywhere. The proper clause
confesses the Father with his only-begotten Son and the Holy Spirit as
one God and one Lord. It then names three distinctions in its confession:
what belongs properly to the Persons, unity in the essence, and equality
in majesty. Its denial of a singularity of one Person protects the real
personal distinctions while confessing one divine substance.

The phrase about receiving the same truth of the Son and Spirit without
a difference of discrimination concerns divine glory. It is followed by
the explicit confession of the Persons' proper distinction. The
historical English's ``without any difference or distinction'' must be
read with that next clause, rather than as cancelling it. Angels and
archangels, cherubim and seraphim praise this God, and the prayer proceeds
to the Sanctus. The human congregation's thanksgiving belongs within
this worship.

Colossians likewise gives thanks to the Father through Christ. That is
a direct point of comparison with the Epistle. The Trinity's unity of
essence is a confession about God; Paul's one-body command concerns
redeemed human members. The moral demand for concord follows Paul's
words about forgiveness and charity.

\subsection{Communion: confidence within believing, forgiving prayer}
Mark 11 places Jesus's instruction after the disciples see the withered
fig tree. He commands faith in God and speaks of the power of believing
prayer. Immediately afterwards he commands people standing in prayer to
forgive anyone against whom they have a grievance (Mark 11:22--26).
The brief Communion chant appoints the promise of reception from this
larger instruction.

Bede attends to the forgiveness command when explaining the askers.
Perfect faith works through love; those who have not reached its fullness
can ask pardon and everlasting life, forgiving their own offenders.
The direction of the prayer and the petitioner's relation to the
neighbour therefore belong within the explanation of confidence.\footnote{Bede,
\work{Exposition of Mark}, on Mark 11:22--26, paragraphs beginning
\textit{Et respondens Jesus} and \textit{Et cum stabitis ad orandum};
retained Latin transcription, raw paragraphs 338--339.}
The second interpretation compares this exposition with Aquinas's account
of the good that prayer obtains.

\subsection{Postcommunion: a received pledge and a requested effect}
The Postcommunion distinguishes two acts. The people have received a
pledge through the mysteries, and they now ask to obtain the effect of
that salvation. Reception has occurred; fulfilment is requested.

This difference between pledge and effect permits an analogy with the
Gospel's time before harvest. Both concern a present whose end is still
awaited. The Postcommunion also specifies what the people desire:
salvation's effect. The Communion's assurance of reception can therefore
be considered alongside an actual saving petition, whose end exceeds
the removal of a present inconvenience.
